\section{The Close of the Year and the Voice of the Judge}

When November strips the trees and shortens the afternoons, the Roman liturgy draws the faithful toward the end of all things. On All Souls' Day the Church vests in black and chants the \textit{Dies irae}, standing in imagination before the throne where every hidden deed comes to light. On the Twenty-fourth and Last Sunday after Pentecost, immediately before the threshold of Advent, the Epistle gives thanks to the Father ``who hath delivered us from the power of darkness and hath translated us into the kingdom of the Son of his love'' (Col 1:13), the Alleluia and Offertory cry out of the depths, \textit{De profundis clamavi ad te, Domine} (Ps 129:1), and the Gospel opens onto the darkened sun, the falling stars, and the sign of the Son of man appearing in heaven (Matt 24:29--30). The Collect of that last Sunday asks God to stir up the wills of His faithful---\textit{Excita, quaesumus, Domine, tuorum fidelium voluntates}---so that they may pursue the fruit of the divine work more eagerly and receive the greater remedies of His mercy.\endnote{1962 \textit{Missale Romanum}, Dominica XXIV et ultima post Pentecosten (Collect, Epistle Col 1:9--14, Alleluia and Offertory Ps 129:1--2, Gospel Matt 24:15--35); Commemoratio omnium fidelium defunctorum, First Mass (\textit{Dies irae}).} At the very close of the Gospel for the dying year stands a single sentence of terrible permanence: ``Heaven and earth shall pass: but my words shall not pass'' (Matt 24:35).

Among the words that shall not pass are the Lord's own warnings about hell. Throughout the New Testament, the Church's teaching on eternal damnation falls chiefly from the lips of Jesus Christ Himself---the same Good Shepherd who sought the lost sheep, wept over Jerusalem, forgave the sinful woman, and opened His arms upon the Cross. No prophet of Israel and no Apostle after Pentecost speaks so plainly or so persistently of eternal ruin as the incarnate Son of God.

In the Sermon on the Mount, Christ warns that whoever calls his brother a fool ``shall be in danger of hell fire''---\textit{gehennae ignis} (Matt 5:22)---and counsels the severest amputation of whatever leads to sin, because ``it is expedient for thee that one of thy members should perish, rather than thy whole body be cast into hell'' (Matt 5:29; cf.\ 5:30; 18:8--9). The Greek word behind His warning, \textit{Geenna}, derives from the Hebrew \textit{Ge Hinnom}, the valley of the sons of Hinnom outside Jerusalem, once defiled by idolatrous sacrifices and later turned into a place of refuse and burning. In the mouth of Christ that valley becomes the name of the final place of perdition, where not only the body perishes: ``And fear ye not them that kill the body, and are not able to kill the soul: but rather fear him that can destroy both soul and body in hell'' (Matt 10:28; cf.\ Luke 12:4--5).

In the Gospel of Mark, the warning takes on a liturgical cadence. Three times in six verses the Lord repeats the same terrible refrain drawn from the closing verse of Isaias (Isa 66:24):

\begin{quote}
And if thy hand scandalize thee, cut it off: it is better for thee to enter into life, maimed, than having two hands to go into hell, into unquenchable fire: Where their worm dieth not, and the fire is not extinguished. And if thy foot scandalize thee, cut it off: it is better for thee to enter lame into life everlasting than having two feet to be cast into the hell of unquenchable fire: Where their worm dieth not, and the fire is not extinguished. And if thy eye scandalize thee, pluck it out: it is better for thee with one eye to enter into the kingdom of God than having two eyes to be cast into the hell of fire: Where their worm dieth not, and the fire is not extinguished. (Mark 9:42--47)\endnote{Verse numbering throughout follows the Clementine Vulgate and Challoner Douay-Rheims Bible, in which the triple refrain stands at Mark 9:43, 45, and 47.}
\end{quote}

St.\ Augustine marvels at this triple repetition in \textit{The City of God}: ``He did not shrink from using the same words three times over in one passage. And who is not terrified by this repetition, and by the threat of that punishment uttered so vehemently by the lips of the Lord Himself?''\endnote{St.\ Augustine, \textit{The City of God} XXI.9, trans.\ Marcus Dods (Edinburgh: T.\ \& T.\ Clark, 1871).} Hand, foot, and eye represent whatever a man clings to as though it were his own flesh; yet no earthly attachment, however dear, is worth keeping at the price of an inextinguishable ruin.

Alongside \textit{Gehenna} and unquenchable fire, our Lord multiplies images so that the mind cannot reduce the reality to a single figure. He speaks of the ``furnace of fire'' into which the angels will cast the cockle and all workers of iniquity at the consummation of the world, where ``there shall be weeping and gnashing of teeth'' (Matt 13:40--42, 49--50). He speaks three times in Matthew's Gospel of ``the exterior darkness''---\textit{tenebras exteriores}---into which the faithless children of the kingdom, the guest without a wedding garment, and the unprofitable servant are cast out (Matt 8:12; 22:13; 25:30). He speaks of the door shut once the bridegroom enters the marriage feast, leaving the foolish virgins outside knocking in vain to hear ``Amen I say to you, I know you not'' (Matt 25:10--12; cf.\ Luke 13:25--27). In the account of the rich man and Lazarus, He shows the dead rich man buried in \textit{infernus}, crying from the flame for a single drop of water on the tip of Lazarus's finger, and hearing Abraham declare that between the blessed and the tormented ``there is fixed a great chaos: so that they who would pass from hence to you cannot, nor from thence come hither'' (Luke 16:22--26). Against the blasphemy of the Holy Ghost---understood by the Church as final impenitence, the willful refusal of forgiveness---He pronounces that the offender ``shall never have forgiveness, but shall be guilty of an everlasting sin'' (Mark 3:29; cf.\ Matt 12:31--32).

All these streams converge in the twenty-fifth chapter of Matthew, where the Son of man comes in majesty with all His angels, gathers all nations, separates the sheep from the goats, and speaks the final sentence to those on His left: ``Depart from me, you cursed, into everlasting fire, which was prepared for the devil and his angels. \dots\ And these shall go into everlasting punishment: but the just, into life everlasting'' (Matt 25:41, 46). The Fathers drew attention to the asymmetry between the two verdicts: the kingdom was ``prepared for you from the foundation of the world'' (Matt 25:34), whereas the everlasting fire was ``prepared for the devil and his angels'' (Matt 25:41). Man was created for the kingdom, not for the pit; he enters that fire only by joining himself to the rebellion of the fallen spirits.

The apostolic writings speak in complete harmony with their Master. St.\ Paul writes to the Thessalonians that when the Lord Jesus is revealed from heaven in a flame of fire, those who know not God and obey not the Gospel ``shall suffer eternal punishment in destruction, from the face of the Lord and from the glory of his power'' (2 Thess 1:8--9). The Epistle to the Hebrews warns those who trample underfoot the Son of God and affront the Spirit of grace that there awaits ``a certain dreadful expectation of judgment, and the rage of a fire which shall consume the adversaries,'' adding, ``It is a fearful thing to fall into the hands of the living God'' (Heb 10:27, 29, 31). St.\ Peter recalls that God spared not the angels that sinned, ``but delivered them, drawn down by infernal ropes to the lower hell, unto torments, to be reserved unto judgment'' (2 Pet 2:4), and St.\ Jude names the ``everlasting chains'' under darkness and the ``storm of darkness'' reserved forever (Jude 1:6, 13). In the Apocalypse, the smoke of the torments of those who adore the beast ascends ``for ever and ever'' (Apoc 14:11), the devil, the beast, and the false prophet are cast into the pool of fire and brimstone to be ``tormented day and night for ever and ever'' (Apoc 20:9--10), and whosoever is not found written in the book of life is cast into that same pool of fire: ``This is the second death'' (Apoc 20:14--15; cf.\ 21:8).

The \textit{Roman Catechism} supplies a necessary lexical distinction when explaining the fifth article of the Apostles' Creed, ``He descended into hell.''\endnote{\textit{Catechism of the Council of Trent}, Part I, Creed Article V, trans.\ J.\ Donovan (Baltimore: Lucas Brothers, 1829), pp.\ 50--51.} In Latin Scripture and the Creed, the word \textit{infernus} or \textit{inferi} (answering to Hebrew \textit{Sheol} and Greek \textit{Hades}) designates in a broad sense the hidden abodes of souls not yet admitted to the beatific vision. The Catechism distinguishes three such abodes: the limbo of the Fathers, where the just who died before Christ awaited redemption in peaceful repose without pain and whom Christ liberated as Victor by His descent; purgatory, where the souls of the just are cleansed by temporary punishment before entering heaven; and ``that most loathsome and dark prison in which the souls of the damned are buried with the unclean spirits, in eternal and inextinguishable fire.'' This last, the Catechism observes, ``is called Gehenna, the bottomless pit, and, strictly speaking, means hell.'' It is of this Gehenna---the second death, eternal exile from the face of God---that the Church speaks when she places hell third among the Four Last Things.
